\documentclass[../../main.tex]{subfiles}

\graphicspath{{../../image/}{../../figure/}}

\begin{document}

\section{双曲平衡律方程组的例子}\label{section:ch03-例子}

在连续介质物理中出现的双曲平衡律方程组不计其数，这里只给出一小批样本。它们将作为灯塔，指引一般理论的发展。

\subsection{标量平衡律}\label{subsection:ch03-标量平衡律}

单个平衡律（$n=1$）
\begin{align}
\partial_{t}u(x,t)+\operatorname{div}G(u(x,t),x,t)=\varpi(u(x,t),x,t)
\label{align:ch03-标量平衡律}
\end{align}
总是双曲的。任意函数 $\eta(u,x,t)$ 都可作为熵，相应的熵流与熵产生由下式算出：
\begin{align}
Q=\int^{u}\frac{\partial\eta}{\partial u}\frac{\partial G}{\partial u}\,\mathrm{d}u,
\label{align:ch03-标量熵流}
\end{align}
\begin{align}
h=\sum_{\alpha=1}^{m}\left[\frac{\partial\eta}{\partial u}\frac{\partial G_{\alpha}}{\partial x_{\alpha}}-\frac{\partial Q_{\alpha}}{\partial x_{\alpha}}\right]+\varpi\frac{\partial\eta}{\partial u}+\frac{\partial\eta}{\partial t}.
\label{align:ch03-标量熵产生}
\end{align}

\eqref{align:ch03-标量平衡律} 以及相应的齐次标量守恒律，特别是它们在一维空间（$m=1$）下的形式，将被广泛用作发展一般方程组理论的模型。

\subsection{无热导热弹性介质}\label{subsection:ch03-无热导热弹性}

热弹性介质的理论已在第~II 章中讨论。这里我们采用参考构型（Lagrange 描述），于是场是 $(x,t)$ 的函数。为与本章记号一致，我们用 $\partial_{t}$ 表示物质时间导数（代替第~II 章中的上点记号），用 $\partial_{\alpha}$ 表示关于 $x$ 的 $\alpha$ 分量 $x_{\alpha}$ 的偏导数。为确定起见，设物理空间维数 $m=3$。我们还采用标准求和约定：重复指标在 $1,2,3$ 范围内求和。

本构方程见\cref{section:ch02-热弹性}%

。由于不再有混淆之虞，我们可以简化记号，删去本构函数符号上的"帽"。同样为简单起见，设介质是均匀的，参考密度 $\rho_{0}=1$。

如第~II 章所解释，一个热力学过程由运动 $\chi$ 与熵场 $s$ 确定。为把平衡律场方程写成\eqref{align:ch03-平衡律一般形式} 形式的一阶方程组，我们通过其导数\eqref{align:ch02-2-1-1}%

、\eqref{align:ch02-2-1-2}%

来监控 $\chi$，从而取状态向量 $U=(F,v,s)$，取值于 $\mathbb{R}^{13}$。在这种情形下，必须在线动量平衡\eqref{align:ch02-2-3-4}%

与能量平衡\eqref{align:ch02-2-3-10}%

之外，再补上协调条件\eqref{align:ch02-2-1-8}$_{1}$%

。于是平衡律方程组为
\begin{align}
\begin{cases}
\partial_{t}F_{i\alpha}-\partial_{\alpha}v_{i}=0, & i,\alpha=1,2,3\\
\partial_{t}v_{i}-\partial_{\alpha}S_{i\alpha}(F,s)=b_{i}, & i=1,2,3\\
\partial_{t}\left[\varepsilon(F,s)+\frac{1}{2}|v|^{2}\right]-\partial_{\alpha}\left[v_{i}S_{i\alpha}(F,s)\right]=b_{i}v_{i}+r, &
\end{cases}
\label{align:ch03-热弹性系统}
\end{align}
其中（见\eqref{align:ch02-2-5-3}%

）
\begin{align}
S_{i\alpha}(F,s)=\frac{\partial\varepsilon(F,s)}{\partial F_{i\alpha}},\qquad \theta(F,s)=\frac{\partial\varepsilon(F,s)}{\partial s}.
\label{align:ch03-热弹性本构}
\end{align}

经过冗长的计算可验证：当且仅当状态空间某区域内的每一 $(F,s)$ 满足
\begin{align}
\frac{\partial\varepsilon(F,s)}{\partial s}>0,
\label{align:ch03-温度正}
\end{align}
\begin{align}
\frac{\partial^{2}\varepsilon(F,s)}{\partial F_{i\alpha}\partial F_{j\beta}}\nu_{\alpha}\nu_{\beta}\xi_{i}\xi_{j}>0,\qquad \text{对一切 }\nu,\xi\in\mathbb{S}^{2}
\label{align:ch03-Legendre-Hadamard}
\end{align}
时，方程组\eqref{align:ch03-热弹性系统} 在该区域上是双曲的。

由\eqref{align:ch03-热弹性本构}$_{2}$，条件\eqref{align:ch03-温度正} 无非是说绝对温度必须为正。\eqref{align:ch03-Legendre-Hadamard} 称为\emph{Legendre--Hadamard 条件}，它意味着 $\varepsilon$ 关于 $F$ 是\emph{秩一凸}的，即沿任何秩一方向 $\xi^{\top}\nu$ 都凸。\eqref{align:ch03-Legendre-Hadamard} 的另一种等价表述是：对任意单位向量 $\nu$，由
\begin{align}
N_{ij}(\nu;F,s)=\frac{\partial^{2}\varepsilon(F,s)}{\partial F_{i\alpha}\partial F_{j\beta}}\nu_{\alpha}\nu_{\beta},\qquad i,j=1,2,3
\label{align:ch03-声张量}
\end{align}
定义的\emph{声张量} $N(\nu;F,s)$ 正定。事实上，对系统\eqref{align:ch03-热弹性系统}，沿方向 $\nu$ 的特征速度是声张量三个特征值的六个平方根，零的重数为七。

回顾第~II 章，除平衡律方程组\eqref{align:ch03-热弹性系统} 外，热力学可容许过程还应满足 Clausius--Duhem 不等式\eqref{align:ch02-2-3-13}%

，在此处写为
\begin{align}
-\partial_{t}s\leqslant-\frac{r}{\theta(F,s)}.
\label{align:ch03-Clausius-Duhem热弹性}
\end{align}

由\eqref{align:ch03-热弹性本构}，\eqref{align:ch03-热弹性系统} 的任一经典解都恒等地把\eqref{align:ch03-Clausius-Duhem热弹性} 作为等式满足。\footnote{因此，对经典解而言，用等式\eqref{align:ch03-Clausius-Duhem热弹性} 代替\eqref{align:ch03-热弹性系统} 中的第三个方程是方便的。特别地，若 $r\equiv0$，熵 $s$ 沿质点轨迹保持常数，于是只需解前两个方程即可确定 $F$ 与 $v$。} 因而，按\cref{section:ch03-熵流对}%

的术语，$-s$ 是方程组\eqref{align:ch03-热弹性系统} 的一个熵，相应熵流为零。\footnote{把 $-s$ 而非物理熵 $s$ 本身识别为"熵"似乎有些奇怪。采用这一约定是因为：处理关于时间单调不增的解泛函更为方便。}\eqref{align:ch03-热弹性系统} 的弱解一般不再满足\eqref{align:ch03-Clausius-Duhem热弹性}。因此\eqref{align:ch03-Clausius-Duhem热弹性} 的作用是剔除不合物理要求的弱解。把一个伴随平衡律从经典解下的等式延拓为弱解下的不等式，这一做法在双曲平衡律方程组的一般理论中起着关键作用。

还应注意：$(F,v,s)$ 是\eqref{align:ch03-热弹性系统} 的解，这一点本身并不使之与弹性动力学相关，除非 $F$ 是变形梯度\eqref{align:ch02-2-1-2}%

，即等价地
\begin{align}
\partial_{\beta}F_{i\alpha}-\partial_{\alpha}F_{i\beta}=0,\qquad i=1,2,3,\quad \alpha,\beta=1,2,3.
\label{align:ch03-协调方程}
\end{align}

在这种情形下，如\cref{section:ch02-连续介质热力学平衡律}%

所示，$(F,v)$ 还将满足运动学守恒律\eqref{align:ch02-2-2-15}%

与\eqref{align:ch02-2-2-16}%

，即
\begin{align}
\partial_{t}F_{\gamma k}^{*}-\partial_{\alpha}\left(\varepsilon_{ijk}\varepsilon_{\alpha\beta\gamma}v_{i}F_{j\beta}\right)=0,\qquad k=1,2,3,\quad \gamma=1,2,3,
\label{align:ch03-余因子守恒}
\end{align}
\begin{align}
\partial_{t}(\det F)-\partial_{\alpha}\left(v_{i}F_{\alpha i}^{*}\right)=0.
\label{align:ch03-行列式守恒}
\end{align}

回顾可知，即使 $F,v$ 仅仅属于 $L^{\infty}$，\eqref{align:ch03-余因子守恒} 与\eqref{align:ch03-行列式守恒} 仍然成立。

沿方向 $\nu\in\mathbb{S}^{2}$ 以速度 $\sigma$ 传播的激波阵面，其 Rankine--Hugoniot 跳跃条件\eqref{align:ch03-RH一般形式} 在此处取为
\begin{align}
\begin{cases}
-\sigma[\![F_{i\alpha}]\!]=[\![v_{i}]\!]\nu_{\alpha}, & i,\alpha=1,2,3\\
-\sigma[\![v_{i}]\!]=[\![S_{i\alpha}(F,s)]\!]\nu_{\alpha}, & i=1,2,3\\
-\sigma[\![\varepsilon(F,s)+\frac{1}{2}|v|^{2}]\!]=[\![v_{i}S_{i\alpha}(F,s)]\!]\nu_{\alpha}, &
\end{cases}
\label{align:ch03-热弹性RH}
\end{align}
其中双方括号 $[\![\,\cdot\,]\!]$ 表示所括量跨过激波的跃度。把\eqref{align:ch03-热弹性RH} 中三个方程联立，可消去速度，得到
\begin{align}
-\sigma\left\{[\![\varepsilon]\!]-\operatorname{tr}\left(\frac{1}{2}(S_{+}+S_{-})^{\top}[\![F]\!]\right)\right\}=0.
\label{align:ch03-热弹性RH消速度}
\end{align}

与物理上有意义的解相关联的任何激波，还必须满足由\eqref{align:ch03-协调方程} 诱导出的跳跃条件
\begin{align}
[\![F_{i\alpha}]\!]\nu_{\beta}=[\![F_{i\beta}]\!]\nu_{\alpha},\qquad i=1,2,3,\quad \alpha,\beta=1,2,3,
\label{align:ch03-协调跳跃}
\end{align}
或等价地，存在向量 $w\in\mathbb{R}^{3}$ 使
\begin{align}
[\![F_{i\alpha}]\!]=w_{i}\nu_{\alpha},\qquad i=1,2,3,\quad \alpha=1,2,3.
\label{align:ch03-协调跳跃分解}
\end{align}

由\eqref{align:ch03-热弹性RH}$_{1}$，当激波速度 $\sigma\neq0$ 时这一条件自动成立，且 $w=\sigma^{-1}[\![v]\!]$。然而\eqref{align:ch03-协调跳跃分解} 排除了任何速度 $\sigma=0$ 的等熵激波。事实上，对连接 $(F_{-},s,v_{-})$ 与 $(F_{+},s,v_{+})$ 的这种激波，\eqref{align:ch03-热弹性RH}$_{2}$ 与\eqref{align:ch03-热弹性本构}、\eqref{align:ch03-声张量}、\eqref{align:ch03-协调跳跃分解} 联立将推出 $\bar{N}w=0$，其中
\begin{align}
\bar{N}=\int_{0}^{1}N\big(\nu,(1-\tau)F_{-}+\tau F_{+},s\big)\,\mathrm{d}\tau.
\label{align:ch03-积分声张量}
\end{align}

由于 $\bar{N}$ 正定，这给出 $w=0$，因而这种激波不可能存在。另一方面，确实存在与\eqref{align:ch03-协调跳跃分解} 相容的静止非等熵激波。

由\eqref{align:ch03-Clausius-Duhem热弹性} 诱导出的\eqref{align:ch01-BV-激波不等式}%

不等式取为
\begin{align}
\sigma[\![s]\!]\leqslant0,
\label{align:ch03-熵增条件热弹性}
\end{align}
它意味着每当物质质点跨过非静止激波时，其物理熵增加。

\subsection{无热导热弹性介质的等熵运动}\label{subsection:ch03-等熵热弹性}

等熵过程的物理背景已在\cref{section:ch02-热弹性}%

中讨论。特别地，如前所述，在无热源 $r=0$ 时，任何初始为等熵的热弹性过程只要保持光滑就始终是等熵的。此外，常熵假设即使对弱解也常常是令人满意的近似。熵被固定为常值 $\bar{s}$，为简单起见从记号中删去。状态向量约化为 $U=(F,v)$，取值于 $\mathbb{R}^{12}$。平衡律方程组由\eqref{align:ch03-热弹性系统} 删去能量平衡得到：
\begin{align}
\begin{cases}
\partial_{t}F_{i\alpha}-\partial_{\alpha}v_{i}=0, & i,\alpha=1,2,3\\
\partial_{t}v_{i}-\partial_{\alpha}S_{i\alpha}(F)=b_{i}, & i=1,2,3
\end{cases}
\label{align:ch03-等熵热弹性系统}
\end{align}
并且仍有
\begin{align}
S_{i\alpha}(F)=\frac{\partial\varepsilon(F)}{\partial F_{i\alpha}},\qquad i,\alpha=1,2,3.
\label{align:ch03-等熵热弹性本构}
\end{align}

方程组\eqref{align:ch03-等熵热弹性系统} 当 $\varepsilon$ 秩一凸（即\eqref{align:ch03-Legendre-Hadamard} 在 $s=\bar{s}$ 处成立）时双曲。沿方向 $\nu\in\mathbb{S}^{2}$ 的特征速度是声张量\eqref{align:ch03-声张量} 在 $s=\bar{s}$ 处三个特征值的六个平方根，零的重数为六。

如\cref{section:ch02-热弹性}%

所解释，除\eqref{align:ch03-等熵热弹性系统} 外，热力学可容许的等熵运动还应满足不等式\eqref{align:ch02-2-5-32}%

，在当前记号下写为
\begin{align}
\partial_{t}\left[\varepsilon(F)+\frac{1}{2}|v|^{2}\right]-\partial_{\alpha}\left[v_{i}S_{i\alpha}(F)\right]\leqslant b_{i}v_{i}.
\label{align:ch03-等熵能量不等式}
\end{align}

由\eqref{align:ch03-等熵热弹性本构}，\eqref{align:ch03-等熵热弹性系统} 的任一经典解都恒等地把\eqref{align:ch03-等熵能量不等式} 作为等式满足。因此按\cref{section:ch03-熵流对}%

术语，$\eta=\varepsilon(F)+\frac{1}{2}|v|^{2}$ 是\eqref{align:ch03-等熵热弹性系统} 的熵。注意\eqref{align:ch03-等熵热弹性系统} 已取正则形式\eqref{align:ch03-正则形式}，且 $D\eta=(S(F),v)$。于是如\cref{section:ch03-熵流对}%

所示，当内能 $\varepsilon(F)$ 一致凸时，把状态向量从 $U=(F,v)$ 改为 $U^{*}=(S,v)$ 将使方程组\eqref{align:ch03-等熵热弹性系统} 成为对称双曲。但应注意：即使 $\varepsilon(F)$ 在状态空间的某一部分（尤其是在其极小点附近）可能凸，它也不可能整体凸，除非它是二次的——而此时\eqref{align:ch03-等熵热弹性系统} 是线性的。这是物质标架无差异原理的推论：该原理要求 $\varepsilon(F)$ 在刚性旋转下不变。

\eqref{align:ch03-等熵热弹性系统} 的弱解一般不再恒等地满足\eqref{align:ch03-等熵能量不等式}。于是我们再次遇到这种情形：伴随平衡律从经典解下的等式延拓为弱解下的不等式，作为弱解的可容许性条件。

与弹性动力学相关的\eqref{align:ch03-等熵热弹性系统} 经典或弱解还应满足\eqref{align:ch03-协调方程}，从而满足运动学守恒律\eqref{align:ch03-余因子守恒} 与\eqref{align:ch03-行列式守恒}。

沿方向 $\nu\in\mathbb{S}^{2}$ 以速度 $\sigma$ 传播的激波阵面，其 Rankine--Hugoniot 跳跃条件取为
\begin{align}
\begin{cases}
-\sigma[\![F_{i\alpha}]\!]=[\![v_{i}]\!]\nu_{\alpha}, & i,\alpha=1,2,3\\
-\sigma[\![v_{i}]\!]=[\![S_{i\alpha}(F)]\!]\nu_{\alpha}, & i=1,2,3.
\end{cases}
\label{align:ch03-等熵热弹性RH}
\end{align}

与物理上有意义的解相关联的激波还应满足\eqref{align:ch03-热弹性RH消速度} 与\eqref{align:ch03-协调跳跃}，如前所述，这排除了所有静止激波。

由\eqref{align:ch03-等熵能量不等式} 诱导出的\eqref{align:ch01-BV-激波不等式}%

不等式为
\begin{align}
-\sigma[\![\varepsilon(F)+\frac{1}{2}|v|^{2}]\!]-[\![v_{i}S_{i\alpha}(F)]\!]\nu_{\alpha}\leqslant0,
\label{align:ch03-等熵熵不等式}
\end{align}
结合\eqref{align:ch03-等熵热弹性RH} 可化为
\begin{align}
-\sigma\left\{[\![\varepsilon]\!]-\operatorname{tr}\left(\frac{1}{2}(S_{+}+S_{-})^{\top}[\![F]\!]\right)\right\}\leqslant0,
\label{align:ch03-等熵RH熵条件}
\end{align}
即\eqref{align:ch03-热弹性RH消速度} 的不含速度的类比。

从\eqref{align:ch03-热弹性系统} 到\eqref{align:ch03-等熵热弹性系统} 的过渡，提供了连续介质物理中常用的一种\emph{截断}过程的范例：删去若干方程，同时按比例相应缩减状态向量的维数，其规则已在\cref{section:ch01-伴随平衡律}%

中给出。事实上，截断方程组\eqref{align:ch03-等熵热弹性系统} 的伴随平衡律\eqref{align:ch03-等熵能量不等式}，可由原方程组\eqref{align:ch03-热弹性系统} 的伴随平衡律\eqref{align:ch03-Clausius-Duhem热弹性} 按\eqref{align:ch01-truncated-companion}%

的配方导出。回顾：在正则截断中，删去任一方程都应同时冻结使方程组对称的特殊状态向量的相应分量。例如，可以删去最后三个方程中的第 $i$ 个、同时冻结速度的第 $i$ 分量 $v_{i}$；或者删去前九个方程中的 $(i,\alpha)$ 个、同时冻结 Piola--Kirchhoff 应力的 $(i,\alpha)$ 分量 $S_{i\alpha}(F)$。

如\cref{section:ch02-热弹性}%

所解释，热弹性材料等温过程的平衡律，可在\eqref{align:ch03-等熵热弹性系统}、\eqref{align:ch03-等熵热弹性本构} 与\eqref{align:ch03-等熵能量不等式} 中把常熵下的内能 $\varepsilon(F)$ 替换为常温下的 Helmholtz 自由能 $\psi(F)$ 而得到。

\subsection{带松弛的等熵运动}\label{subsection:ch03-带松弛等熵}

我们考虑\cref{section:ch02-松弛}%

中所述材料的等熵运动，为简单起见设参考密度 $\rho_{0}=1$、体积力 $b=0$。状态向量为 $U=(F,v,\Sigma)$，取值于 $\mathbb{R}^{21}$。平衡律方程组由协调方程\eqref{align:ch02-2-1-8}$_{1}$%

、线动量平衡\eqref{align:ch02-2-3-4}%

以及内变量 $\Sigma$ 的平衡律\eqref{align:ch02-2-8-3}%

组成：
\begin{align}
\begin{cases}
\partial_{t}F_{i\alpha}-\partial_{\alpha}v_{i}=0, & i,\alpha=1,2,3\\
\partial_{t}v_{i}-\partial_{\alpha}\left[P_{i\alpha}(F)+\Sigma_{i\alpha}\right]=0, & i=1,2,3\\
\partial_{t}\Sigma_{i\alpha}=\frac{1}{\tau}\left[\Pi_{i\alpha}(\Sigma)-F_{i\alpha}\right], & i,\alpha=1,2,3.
\end{cases}
\label{align:ch03-松弛系统}
\end{align}
此外，
\begin{align}
P_{i\alpha}(F)=\frac{\partial\sigma(F)}{\partial F_{i\alpha}},\qquad \Pi_{i\alpha}(\Sigma)=-\frac{\partial h(\Sigma)}{\partial\Sigma_{i\alpha}}.
\label{align:ch03-松弛本构}
\end{align}

除\eqref{align:ch03-松弛系统} 外，热力学可容许（等熵）运动还应满足熵不等式\eqref{align:ch02-2-5-32}%

，在此处写为
\begin{align}
\partial_{t}\left[\sigma(F)+\operatorname{tr}(\Sigma F^{\top})+h(\Sigma)+\frac{1}{2}|v|^{2}\right]-\partial_{\alpha}\left[v_{i}P_{i\alpha}(F)+v_{i}\Sigma_{i\alpha}\right]\leqslant0,
\label{align:ch03-松弛熵不等式}
\end{align}
于是按\cref{section:ch03-熵流对}%

术语，$\sigma(F)+\operatorname{tr}(\Sigma F^{\top})+h(\Sigma)+\frac{1}{2}|v|^{2}$ 是\eqref{align:ch03-松弛系统} 的熵。

方程组\eqref{align:ch03-松弛系统} 当
\begin{align}
\frac{\partial^{2}\sigma(F)}{\partial F_{i\alpha}\partial F_{j\beta}}\nu_{\alpha}\nu_{\beta}\xi_{i}\xi_{j}+\nu_{\alpha}\nu_{\alpha}\xi_{i}\xi_{i}+\frac{\partial^{2}h(\Sigma)}{\partial\Sigma_{i\alpha}\partial\Sigma_{j\beta}}\nu_{\alpha}\nu_{\beta}\zeta_{i}\zeta_{j}>0
\label{align:ch03-松弛双曲条件}
\end{align}
对一切 $\nu\in\mathbb{S}^{2}$ 与 $(\xi,\zeta)^{\top}\in\mathbb{S}^{5}$ 成立时双曲。

\subsection{无热导热弹性流体}\label{subsection:ch03-热弹性流体}

平衡律方程组\eqref{align:ch03-热弹性系统} 支配所有热弹性介质（特别是热弹性流体）的绝热热力学过程。但在后一情形下，采用空间描述（Euler 描述）更为有利。原因在于：如\cref{section:ch02-热弹性}%

所示，热弹性流体的内能、温度与 Cauchy 应力只通过密度 $\rho$ 依赖于变形梯度 $F$。于是可以弃用 $F$，而用状态向量 $U=(\rho,v,s)$ 描述介质，取值于（维数小得多的）$\mathbb{R}^{5}$。

场现在是 $(\chi,t)$ 的函数。但为与本章记号一致，我们把符号 $\chi$ 换成 $x$。并且用 $\partial_{t}$（而不用第~II 章中的 $t$ 下标）表示关于 $t$ 的偏导数。

所考虑的平衡律为质量平衡\eqref{align:ch02-2-3-2}%

、线动量平衡\eqref{align:ch02-2-3-5}%

与能量平衡\eqref{align:ch02-2-3-11}%

。本构关系为\eqref{align:ch02-2-5-12}%

（其中 $\tilde{q}\equiv0$）、\eqref{align:ch02-2-5-13}%

与\eqref{align:ch02-2-5-14}%

。为简化记号，我们删去"波浪号"，把 $\tilde{\varepsilon}(\rho,s)$ 写为 $\varepsilon(\rho,s)$。于是平衡律方程组取为
\begin{align}
\begin{cases}
\partial_{t}\rho+\operatorname{div}(\rho v^{\top})=0\\
\partial_{t}(\rho v)+\operatorname{div}(\rho vv^{\top})+\operatorname{grad}p(\rho,s)=\rho b,\\
\partial_{t}\left[\rho\varepsilon(\rho,s)+\frac{1}{2}\rho|v|^{2}\right]+\operatorname{div}\left[\left(\rho\varepsilon(\rho,s)+\frac{1}{2}\rho|v|^{2}+p(\rho,s)\right)v^{\top}\right]=\rho b\cdot v+\rho r,
\end{cases}
\label{align:ch03-Euler}
\end{align}
其中
\begin{align}
p(\rho,s)=\rho^{2}\varepsilon_{\rho}(\rho,s),\qquad \theta(\rho,s)=\varepsilon_{s}(\rho,s).
\label{align:ch03-Euler本构}
\end{align}

方程组\eqref{align:ch03-Euler} 当
\begin{align}
\varepsilon_{s}(\rho,s)>0,\qquad p_{\rho}(\rho,s)>0
\label{align:ch03-Euler双曲条件}
\end{align}
时双曲。沿方向 $\nu\in\mathbb{S}^{2}$ 的特征速度为：重数三的 $v\cdot\nu$，以及重数一的 $v\cdot\nu\pm\sqrt{p_{\rho}(\rho,s)}$。量 $c(\rho,s)=\sqrt{p_{\rho}(\rho,s)}$ 表示随流体流动的观察者所感知的弱波传播速度，称为\emph{声速}。

除\eqref{align:ch03-Euler} 外，热力学可容许过程还应满足 Clausius--Duhem 不等式\eqref{align:ch02-2-3-14}%

，在此处简化为
\begin{align}
\partial_{t}(-\rho s)+\operatorname{div}(-\rho s v^{\top})\leqslant-\rho\frac{r}{\theta(\rho,s)}.
\label{align:ch03-Euler熵不等式}
\end{align}

当过程光滑时，由\eqref{align:ch03-Euler} 与\eqref{align:ch03-Euler本构} 可知\eqref{align:ch03-Euler熵不等式} 恒等地作为等式成立。\footnote{因此对光滑解而言，用较简单的等式\eqref{align:ch03-Euler熵不等式} 代替\eqref{align:ch03-Euler} 的第三个方程常常是方便的。注意：当 $\rho,v,s$ 仅仅 Lipschitz 连续时，光滑性要求即已满足，这允许存在带弱波的流。} 因而 $\eta=-\rho s$ 是方程组\eqref{align:ch03-Euler} 的熵，相应熵流为 $-\rho s v^{\top}$。我们再次看到：伴随平衡律从经典解下的等式延拓为弱解下的不等式，作为弱解物理可容许性的检验。

把状态变量从 $(\rho,v,s)$ 改为 $(\rho,m,E)$，其中 $m=\rho v$ 是动量密度、$E=\rho\varepsilon+\frac{1}{2}\rho|v|^{2}$ 是能量密度，即把\eqref{align:ch03-Euler} 化为其正则形式。一个冗长但常规的计算表明，由\eqref{align:ch03-Euler双曲条件}，熵 $\eta=-\rho s$ 是 $(\rho,m,E)$ 的凸函数。

沿方向 $\nu\in\mathbb{S}^{2}$ 以速度 $\sigma$ 传播的激波阵面，其 Rankine--Hugoniot 跳跃条件如下：
\begin{align}
\begin{cases}
[\![\rho(v\cdot\nu-\sigma)]\!]=0\\
[\![\rho(v\cdot\nu-\sigma)v+p(\rho,s)\nu]\!]=0\\
[\![\rho(v\cdot\nu-\sigma)\left(\varepsilon(\rho,s)+\frac{1}{2}|v|^{2}\right)+p(\rho,s)v\cdot\nu]\!]=0.
\end{cases}
\label{align:ch03-Euler-RH}
\end{align}

联立这些方程，可导出\eqref{align:ch03-热弹性RH消速度} 的类比：
\begin{align}
\rho_{\pm}(v_{\pm}\cdot\nu-\sigma)\left\{[\![\varepsilon]\!]+\frac{1}{2}(p_{+}+p_{-})[\![\rho^{-1}]\!]\right\}=0.
\label{align:ch03-Euler-RH消速度}
\end{align}

注意\eqref{align:ch03-Euler-RH} 允许 $v_{\pm}\cdot\nu=\sigma$ 且 $[\![p]\!]=0$ 的激波。这些波以特征速度传播，称为\emph{接触间断}或\emph{涡层}。流体质点可在涡层两侧以不同速度滑动，但不能穿过它。此外，\eqref{align:ch03-Euler-RH} 还支持 $v_{\pm}\cdot\nu\neq\sigma$ 的激波，流体质点轨迹穿过这种激波。\eqref{align:ch03-Euler-RH消速度} 花括号中的量沿这些激波为零。

由\eqref{align:ch03-Euler熵不等式} 诱导出的\eqref{align:ch01-BV-激波不等式}%

不等式取为
\begin{align}
\rho_{\pm}(v_{\pm}\cdot\nu-\sigma)[\![s]\!]\geqslant0,
\label{align:ch03-Euler熵增}
\end{align}
它意味着当流体质点跨过激波时物理熵增加。

\subsection{热弹性流体的等熵流}\label{subsection:ch03-等熵流}

如上所述，在无热源 $r=0$ 时，初始等熵的热弹性流只要保持光滑就始终等熵。此外，即使在激波形成之后，只要激波振幅很小，常熵假设仍然令人满意。在等熵流中，熵被固定为常值并从记号中删去。状态向量约化为 $U=(\rho,v)$，取值于 $\mathbb{R}^{4}$。平衡律方程组由\eqref{align:ch03-Euler} 删去能量平衡得到：
\begin{align}
\begin{cases}
\partial_{t}\rho+\operatorname{div}(\rho v^{\top})=0\\
\partial_{t}(\rho v)+\operatorname{div}(\rho vv^{\top})+\operatorname{grad}p(\rho)=\rho b,
\end{cases}
\label{align:ch03-等熵Euler}
\end{align}
其中
\begin{align}
p(\rho)=\rho^{2}\varepsilon'(\rho).
\label{align:ch03-等熵Euler本构}
\end{align}

方程组\eqref{align:ch03-等熵Euler} 当
\begin{align}
p'(\rho)>0
\label{align:ch03-等熵Euler双曲条件}
\end{align}
时双曲。沿方向 $\nu\in\mathbb{S}^{2}$ 的特征速度为：重数二的 $v\cdot\nu$，以及重数一的 $v\cdot\nu\pm\sqrt{p'(\rho)}$。量 $c(\rho)=\sqrt{p'(\rho)}$ 是声速。特别地，对等熵理想气体\eqref{align:ch02-2-5-31}%

，只要 $\rho>0$，\eqref{align:ch03-等熵Euler双曲条件} 即成立。

热力学可容许的等熵运动还应满足不等式\eqref{align:ch02-2-5-33}%

，在此处简化为
\begin{align}
\partial_{t}\left[\rho\varepsilon(\rho)+\frac{1}{2}\rho|v|^{2}\right]+\operatorname{div}\left[\left(\rho\varepsilon(\rho)+\frac{1}{2}\rho|v|^{2}+p(\rho)\right)v^{\top}\right]\leqslant\rho b\cdot v.
\label{align:ch03-等熵Euler能量不等式}
\end{align}

应注意，方程组\eqref{align:ch03-等熵Euler} 是按\cref{section:ch01-伴随平衡律}%

所述的正则截断由\eqref{align:ch03-Euler} 得到的；特别地，伴随平衡律\eqref{align:ch03-等熵Euler能量不等式} 可由\eqref{align:ch03-Euler熵不等式} 按\eqref{align:ch01-truncated-companion}%

导出。

至此模式已很熟悉：由\eqref{align:ch03-等熵Euler本构}，\eqref{align:ch03-等熵Euler} 的任一经典解都恒等地把\eqref{align:ch03-等熵Euler能量不等式} 作为等式满足，于是函数 $\eta=\rho\varepsilon(\rho)+\frac{1}{2}\rho|v|^{2}$ 是\eqref{align:ch03-等熵Euler} 的熵。同时，不等式\eqref{align:ch03-等熵Euler能量不等式} 被用来剔除物理上不可容许的弱解。

方程组\eqref{align:ch03-等熵Euler} 把状态变量从 $(\rho,v)$ 改为 $(\rho,m)$（$m=\rho v$ 为动量密度）后即取正则形式。容易看出，由\eqref{align:ch03-等熵Euler本构} 与\eqref{align:ch03-等熵Euler双曲条件}，上述熵 $\eta$ 是 $(\rho,m)$ 的凸函数。

沿方向 $\nu\in\mathbb{S}^{2}$ 以速度 $\sigma$ 传播的激波阵面，其 Rankine--Hugoniot 跳跃条件取为
\begin{align}
\begin{cases}
[\![\rho(v\cdot\nu-\sigma)]\!]=0\\
[\![\rho(v\cdot\nu-\sigma)v+p(\rho)\nu]\!]=0.
\end{cases}
\label{align:ch03-等熵Euler-RH}
\end{align}

与非等熵情形一样，存在 $v_{\pm}\cdot\nu=\sigma$ 且 $[\![p]\!]=0$ 的接触间断（涡层），也存在 $v_{\pm}\cdot\nu\neq\sigma$ 的激波。

由\eqref{align:ch03-等熵Euler能量不等式} 诱导出的\eqref{align:ch01-BV-激波不等式}%

不等式为
\begin{align}
[\![\rho(v\cdot\nu-\sigma)\left(\varepsilon(\rho)+\frac{1}{2}|v|^{2}\right)+p(\rho)v\cdot\nu]\!]\leqslant0,
\label{align:ch03-等熵Euler熵不等式}
\end{align}
由\eqref{align:ch03-等熵Euler-RH} 可写为
\begin{align}
\rho_{\pm}(v_{\pm}\cdot\nu-\sigma)\left\{[\![\varepsilon]\!]+\frac{1}{2}(p_{+}+p_{-})[\![\rho^{-1}]\!]\right\}\leqslant0.
\label{align:ch03-等熵Euler-RH熵条件}
\end{align}

对包括多方气体在内的一大类状态方程，\eqref{align:ch03-等熵Euler-RH熵条件} 意味着：当流体质点跨过激波时，密度增大而法向速度减小。

现在设流是光滑等熵的，且体积力来自势
\begin{align}
b=-\operatorname{grad}g,
\label{align:ch03-体积力势}
\end{align}
并监控\cref{section:ch02-运动学}%

引入的自旋张量 $W$ 的演化。把\eqref{align:ch03-等熵Euler} 中两个方程联立，得
\begin{align}
\partial_{t}v+Lv+\operatorname{grad}[h(\rho)+g]=0,
\label{align:ch03-Euler动量}
\end{align}
其中 $L$ 是速度梯度，$h=\varepsilon+p/\rho$ 是\emph{焓}，其导数 $h'(\rho)=p'(\rho)/\rho$。由\eqref{align:ch02-2-1-9}%

，$Lv=2Wv+L^{\top}v$，且 $L^{\top}v=\operatorname{grad}\left(\frac{1}{2}|v|^{2}\right)$，于是
\begin{align}
\partial_{t}v+2Wv+\operatorname{grad}\left[h(\rho)+\frac{1}{2}|v|^{2}+g\right]=0.
\label{align:ch03-Euler动量2}
\end{align}

对\eqref{align:ch03-Euler动量2} 求导并取反称部分，得
\begin{align}
\partial_{t}W+(\mathrm{d}W)v+WL-L^{\top}W^{\top}=0,
\label{align:ch03-涡量方程}
\end{align}
或，利用\eqref{align:ch02-2-1-4}%

与\eqref{align:ch02-2-1-9}%

，
\begin{align}
\dot{W}+WD+DW=0.
\label{align:ch03-涡量方程2}
\end{align}

于是，当 $W$（从而涡量 $\omega=\operatorname{curl}v$）在时空某一点为零时，它沿占据该位置的质点轨迹处处为零。特别地，若 $t=0$ 时对所有 $x$ 涡量都为零，则只要流保持光滑，它在各处都为零。这种流称为\emph{无旋}或\emph{势流}，因为速度场来自势 $\phi$：
\begin{align}
v=\operatorname{grad}\phi.
\label{align:ch03-势流}
\end{align}

把\eqref{align:ch03-势流} 代入\eqref{align:ch03-Euler动量2}，令 $W=0$，积分所得方程，可推出无旋流满足\emph{Bernoulli 方程}
\begin{align}
\partial_{t}\phi+\frac{1}{2}|\operatorname{grad}\phi|^{2}+h(\rho)+g=0.
\label{align:ch03-Bernoulli}
\end{align}
积分常数已被吸收进 $\partial_{t}\phi$ 项中。

于是无旋流可由 Bernoulli 方程\eqref{align:ch03-Bernoulli} 与连续性方程\eqref{align:ch03-等熵Euler}$_{1}$ 联立确定，后者现在写为
\begin{align}
\partial_{t}\rho+\operatorname{div}\left[\rho(\operatorname{grad}\phi)^{\top}\right]=0.
\label{align:ch03-连续性势流}
\end{align}

上述分析即使在 $\rho,v$ 仅仅 Lipschitz 连续时仍然成立，即无旋性在出现弱波时仍保持；但当跳跃间断形成时一般失效，因为激波会产生涡量。然而，把方程组\eqref{align:ch03-Bernoulli}、\eqref{align:ch03-连续性势流} 用于弱解区域已是常见做法，尽管所得流可能不满足动量平衡方程\eqref{align:ch03-等熵Euler}$_{2}$。这与考虑间断等熵流时 $s=$ 常数、尽管它可能不满足能量平衡方程\eqref{align:ch03-Euler}$_{3}$ 是同一精神。与等熵近似类似，当间断（激波）很弱时误差很小。一个重大限制是：无旋流不能承载涡层。事实上，由于 $\phi$ Lipschitz，其切向导数跨过跳跃间断必连续；而表征涡层的条件 $v_{\pm}\cdot\nu=\sigma$ 将推出 $\phi$ 的法向导数也连续，因而这种跳跃间断不可能存在。

等温流的方程可在\eqref{align:ch03-等熵Euler}、\eqref{align:ch03-等熵Euler本构}、\eqref{align:ch03-等熵Euler能量不等式} 中把常熵下的内能 $\varepsilon(\rho)$ 替换为常温下的 Helmholtz 自由能 $\psi(\rho)$ 而得到。

热弹性流体的等熵流与等温流是\emph{正压流}的例子：压力只依赖于密度。

\subsection{Boltzmann 方程与扩展热力学}\label{subsection:ch03-Boltzmann}

与连续介质物理相对，动理论把物质视为相互作用分子的集合，并用分子密度函数 $f(\xi,x,t)$ 刻画状态，其中 $\xi\in\mathbb{R}^{3}$ 是占据位置 $x\in\mathbb{R}^{3}$、时刻 $t$ 的分子速度。在适用于单原子气体的经典动理论中，$f(\xi,x,t)$ 满足 Boltzmann 方程
\begin{align}
\partial_{t}f+\xi\cdot\operatorname{grad}_{x}f=\mathcal{Q}[f],
\label{align:ch03-Boltzmann}
\end{align}
其中 $\mathcal{Q}$ 是一个复杂的积分算子，刻画分子碰撞对 $f$ 的改变。

通过监控密度 $f$ 的各阶矩
\begin{align}
F_{i_{1}\cdots i_{N}}=\int_{\mathbb{R}^{3}}\xi_{i_{1}}\cdots\xi_{i_{N}}f\,\mathrm{d}\xi,\qquad i_{1},\cdots,i_{N}=1,2,3
\label{align:ch03-矩}
\end{align}
可以建立连续介质方法与动理论之间的形式联系。事实上，这些矩满足一个无穷阶发展方程组
\begin{align}
\begin{cases}
\partial_{t}F+\partial_{j}F_{j}=0\\
\partial_{t}F_{i}+\partial_{j}F_{ij}=0, & i=1,2,3\\
\partial_{t}F_{ij}+\partial_{k}F_{ijk}=P_{ij}, & i,j=1,2,3\\
\partial_{t}F_{ijk}+\partial_{\ell}F_{ijk\ell}=P_{ijk}, & i,j,k=1,2,3\\
\cdots\cdots\cdots\cdots\cdots\\
\partial_{t}F_{i_{1}\cdots i_{N}}+\partial_{m}F_{i_{1}\cdots i_{N}m}=P_{i_{1}\cdots i_{N}}, & i_{1},\ldots,i_{N}=1,2,3.
\end{cases}
\label{align:ch03-矩方程组}
\end{align}

在上述方程以及本节各处，$\partial_{i}=\partial/\partial x_{i}$，并采用求和约定：重复指标在 $1,2,3$ 范围内求和。$P_{i_{1}\cdots i_{N}}$ 项表示 $\xi_{i_{1}}\cdots\xi_{i_{N}}\mathcal{Q}[f]$ 在 $\mathbb{R}^{3}$ 上的积分。由于 $\mathcal{Q}$ 的特殊结构，$P_{ij}$ 的迹 $P_{ii}=0$。

注意\eqref{align:ch03-矩方程组} 中每个方程都可视为连续介质物理意义下的平衡律。在这一解释下，$f$ 的矩同时扮演平衡广延量的密度与通量角色。事实上，每一方程中的通量成为下一方程中的密度。在如下识别下
\begin{align}
F=\rho,
\label{align:ch03-识别F}
\end{align}
\begin{align}
F_{i}=\rho v_{i},\qquad i=1,2,3,
\label{align:ch03-识别Fi}
\end{align}
\begin{align}
F_{ij}=\rho v_{i}v_{j}-T_{ij},\qquad i,j=1,2,3,
\label{align:ch03-识别Fij}
\end{align}
\begin{align}
\frac{1}{2}F_{ii}=\rho\varepsilon+\frac{1}{2}\rho|v|^{2},
\label{align:ch03-识别Fii}
\end{align}
\begin{align}
\frac{1}{2}F_{iik}=\left(\rho\varepsilon+\frac{1}{2}\rho|v|^{2}\right)v_{k}-T_{ki}v_{i}+q_{k},\qquad k=1,2,3,
\label{align:ch03-识别F iik}
\end{align}
\eqref{align:ch03-矩方程组} 中第一个方程给出质量守恒，第二个给出线动量守恒，第三个方程取迹的一半给出粘性导热气体的能量守恒，其中 $\rho$ 为密度、$v$ 为速度、$\varepsilon$ 为内能、$T$ 为 Cauchy 应力、$q$ 为热流。我们把 $T$ 视为压力 $p=-\frac{1}{3}T_{ii}$ 与无迹偏斜应力 $\sigma$ 的叠加 $T=-pI+\sigma$，其中 $\sigma_{ii}=0$。由\eqref{align:ch03-识别Fij} 与\eqref{align:ch03-识别Fii}，
\begin{align}
\rho\varepsilon=\frac{3}{2}p,
\label{align:ch03-理想气体关系}
\end{align}
这与理想气体本构关系\eqref{align:ch02-2-5-20}%

在 $\gamma=5/3$ 时相容。

受上述观察启发，可以通过截断无穷方程组\eqref{align:ch03-矩方程组}、只保留有限个方程来构造完整的连续介质理论层级。但所得方程组并不封闭，因为最高阶矩（出现在最后一两个方程中作为通量）以及右端的产生项仍未确定。按连续介质物理的精神，\emph{扩展热力学} 通过如下假设封闭这些方程组：最高阶矩与产生项通过本构关系与较低阶矩相联系，而本构关系由要求方程组的所有光滑解都恒等地满足一个类似于 Clausius--Duhem 不等式的不等式来确定。这诱导出一个伴随平衡律，使方程组可对称化，从而双曲。本构关系还应满足物质标架无差异原理。

为看该方案如何实际运作，我们构造\eqref{align:ch03-矩方程组} 的一个截断，取状态向量 $U=(\rho,v,p,\sigma,q)$，维数为 13（因 $\sigma$ 对称且无迹）。为此保留\eqref{align:ch03-矩方程组} 前三个方程共 10 个独立标量方程，并从第四个方程通过缩并两个指标再取出 3 个方程。由\eqref{align:ch03-识别F}--\eqref{align:ch03-识别F iik}，并注意 $P,P_{i},P_{ii}$ 都为零，最终得方程组
\begin{align}
\begin{cases}
\partial_{t}\rho+\partial_{j}(\rho v_{j})=0\\
\partial_{t}(\rho v_{i})+\partial_{j}(\rho v_{i}v_{j}+p\delta_{ij}-\sigma_{ij})=0\\
\partial_{t}\left(\rho\varepsilon+\frac{1}{2}\rho|v|^{2}\right)+\partial_{k}\left\{\left(\rho\varepsilon+\frac{1}{2}\rho|v|^{2}+p\right)v_{k}-\sigma_{kj}v_{j}+q_{k}\right\}=0\\
\partial_{t}\left(\rho v_{i}v_{j}-\frac{1}{3}\rho|v|^{2}\delta_{ij}-\sigma_{ij}\right)+\partial_{k}\left(F_{ijk}-\frac{1}{3}F_{\ell\ell k}\delta_{ij}\right)=P_{ij}\\
\partial_{t}\left\{\left(\rho\varepsilon+\frac{1}{2}\rho|v|^{2}+p\right)v_{k}-\sigma_{kj}v_{j}+q_{k}\right\}+\frac{1}{2}\partial_{\ell}F_{jj i\ell}=\frac{1}{2}P_{iik}.
\end{cases}
\label{align:ch03-扩展热力学系统}
\end{align}

该方程组通过假设 $F_{ijk},F_{jjik},P_{ij},P_{iik}$ 都是状态向量 $U=(\rho,v,p,\sigma,q)$ 的函数而封闭；这些函数由要求所有光滑解都恒等地满足一个不等式来确定：
\begin{align}
\partial_{t}\varphi+\partial_{i}\psi_{i}\leqslant0,
\label{align:ch03-扩展热力学熵不等式}
\end{align}
其中 $\varphi$ 与 $\psi_{i}$ 是 $U$ 的（未指定）函数，且 $\varphi(U)$ 凸。经过冗长的计算（见\cref{section:ch03-注记}%

引用的文献），可得虽复杂但相当显式的本构关系：
\begin{align}
F_{ijk}=\rho v_{i}v_{j}v_{k}+\left(pv_{k}+\frac{2}{5}q_{k}\right)\delta_{ij}+\left(pv_{i}+\frac{2}{5}q_{i}\right)\delta_{jk}+\left(pv_{j}+\frac{2}{5}q_{j}\right)\delta_{ik},
\label{align:ch03-Fijk}
\end{align}
\begin{align}
\begin{aligned}
F_{jjik}&=\left(\rho|v|^{2}+7p\right)v_{i}v_{k}+\left(p\delta_{ik}-\sigma_{ik}\right)|v|^{2}-\sigma_{ij}v_{j}v_{k}\\
&\quad-\sigma_{kj}v_{j}v_{i}+\frac{14}{5}\left(q_{i}v_{k}+q_{k}v_{i}\right)+\frac{4}{5}q_{j}v_{j}\delta_{ik}+\frac{p}{\rho}\left(5p\delta_{ik}-7\sigma_{ik}\right),
\end{aligned}
\label{align:ch03-Fjjik}
\end{align}
\begin{align}
P_{ij}=\tau_{0}\sigma_{ij},\qquad P_{iik}=2\tau_{0}\sigma_{ki}v_{i}-\tau_{1}q_{k}.
\label{align:ch03-Pij}
\end{align}

为完整起见，$p,\tau_{0},\tau_{1}$ 必须指定为 $(\rho,\theta)$ 的函数。

按\cref{section:ch01-伴随平衡律}%

记号，使方程组对称的特殊向量 $U^{*}=B(U)$ 具有分量
\begin{align}
U^{*}=\frac{1}{\theta}
\begin{pmatrix}
\frac{5p}{2\rho}-\theta s-\frac{1}{2}|v|^{2}+\frac{1}{2p}\sigma_{ij}v_{i}v_{j}-\frac{\rho}{5p^{2}}q_{i}v_{i}|v|^{2}\\
v_{i}-\frac{1}{p}\sigma_{ij}v_{j}+\frac{\rho}{5p^{2}}\left(|v|^{2}q_{i}+2q_{j}v_{j}v_{i}\right)\\
-1+\frac{2\rho}{3p^{2}}q_{k}v_{k}\\
-\frac{1}{2p}\sigma_{ij}-\frac{\rho}{5p^{2}}\left(v_{i}q_{j}+v_{j}q_{i}-\frac{2}{3}v_{k}q_{k}\delta_{ij}\right)\\
\frac{\rho}{5p^{2}}q_{i}
\end{pmatrix}.
\label{align:ch03-对称变量}
\end{align}

特别地，如\cref{section:ch01-伴随平衡律}%

所解释，通过删去最后两个方程来截断方程组\eqref{align:ch03-扩展热力学系统} 时，应同时"冻结" $U^{*}$ 的最后两个分量，即令 $q=0$ 与 $\sigma=0$。此时\eqref{align:ch03-扩展热力学系统} 前三个方程在 $b=0$、$r=0$ 且 $\rho\varepsilon$ 与 $p$ 由\eqref{align:ch03-理想气体关系} 联系的特殊情形下，约化为方程组\eqref{align:ch03-Euler}。若把 $(\rho,v,p)$ 解释为基本状态变量、$(\sigma,q)$ 解释为内状态变量（如\cref{section:ch02-松弛}%

所述），则\eqref{align:ch03-Euler} 成为\eqref{align:ch03-扩展热力学系统} 的松弛形式。

\subsection{非线性电动力学}\label{subsection:ch03-非线性电动力学}

电磁学是另一类有趣的双曲平衡律方程组的丰富源泉。其底层方程组由 Maxwell 方程组成
\begin{align}
\begin{cases}
\partial_{t}B=-\operatorname{curl}E\\
\partial_{t}D=\operatorname{curl}H-J
\end{cases}
\label{align:ch03-Maxwell}
\end{align}
\begin{align}
\operatorname{div}B^{\top}=0,\qquad \operatorname{div}D^{\top}=\rho
\label{align:ch03-高斯}
\end{align}
在 $\mathbb{R}^{3}$ 上，联系电场 $E$、磁场 $H$、电位移 $D$、磁感应 $B$、电流 $J$ 与电荷 $\rho$。电流与电荷由连续性方程相互联系：
\begin{align}
\partial_{t}\rho+\operatorname{div}J^{\top}=0.
\label{align:ch03-电荷守恒}
\end{align}

本构关系由状态向量 $U=(B,D)$ 确定 $E$ 与 $H$。例如，当介质是均匀导电体、具有线性介电响应、相对惯性系静止时，$D=\varepsilon E$，$B=\mu H$，其中 $\varepsilon$ 是介电常数、$\mu$ 是磁导率。为考虑非线性介电响应与电磁场的交叉耦合，人们假设一般本构方程
\begin{align}
E=\frac{\partial\eta(B,D)}{\partial D},\qquad H=\frac{\partial\eta(B,D)}{\partial B},
\label{align:ch03-本构eta}
\end{align}
或
\begin{align}
D=\frac{\partial h(B,E)}{\partial E},\qquad H=-\frac{\partial h(B,E)}{\partial B},
\label{align:ch03-本构h}
\end{align}
其中 $\eta$ 是\emph{电磁场能量}，$h$ 是\emph{Lagrange 量}。注意 $\eta$ 是 $h$ 关于 $E$ 的 Legendre 变换。

物理上可容许的场还应满足耗散不等式
\begin{align}
\partial_{t}\eta(B,D)+\operatorname{div}Q^{\top}(B,D)\leqslant-J\cdot E,
\label{align:ch03-电磁耗散}
\end{align}
其中
\begin{align}
Q=E\wedge H
\label{align:ch03-Poynting}
\end{align}
是\emph{Poynting 向量}。一个直接的计算表明，\eqref{align:ch03-Maxwell}、\eqref{align:ch03-本构eta} 的光滑解恒等地把\eqref{align:ch03-电磁耗散} 作为等式满足。因此 $(\eta,Q)$ 是方程组\eqref{align:ch03-Maxwell}、\eqref{align:ch03-本构eta} 的熵--熵流对。由于 $D\eta=(H,E)$，由\cref{section:ch03-熵流对}%

的讨论可知：当电磁场能量函数一致凸时，把状态向量从 $U=(B,D)$ 改为 $U^{*}=(H,E)$ 使方程组成为对称双曲。但应注意：即使 $\eta$ 在状态空间的某一部分（尤其是在其极小点附近）可能凸，它也不可能整体凸，除非它是二次的——而此时\eqref{align:ch03-Maxwell} 是线性的。这是 Lagrange 量 $h(B,E)$ 在 Lorentz 变换下不变这一要求的推论，因而 $h(B,E)$ 只能通过标量 $|B|^{2}-|E|^{2}$ 与 $B\cdot E$ 依赖于 $(B,E)$。这反过来又意味着 $\eta$ 在刚性旋转下不变，从而只能通过标量 $|B|^{2},|D|^{2}$ 与 $B\cdot D$ 依赖于 $(B,D)$。于是这种函数不可能整体凸，除非它是二次的。在这方面，电磁学与弹性动力学之间存在显著相似性。Lagrange 量 Lorentz 不变性的另一个重要推论是 $E\wedge H=B\wedge D$。Born--Infeld 本构关系\footnote{它们的设计初衷是：与经典线性理论相反，静止点电荷产生的电磁能量有限。} 提供了一个例子：
\begin{align}
\begin{cases}
E=\dfrac{\partial\eta}{\partial D}=\dfrac{1}{\eta}[D+B\wedge Q], & H=\dfrac{\partial\eta}{\partial B}=\dfrac{1}{\eta}[B-D\wedge Q]\\
\eta=\sqrt{1+|B|^{2}+|D|^{2}+|Q|^{2}}, & Q=D\wedge B=E\wedge H.
\end{cases}
\label{align:ch03-BornInfeld}
\end{align}

这里 $\eta$ 不是整体凸的，但所得方程组\eqref{align:ch03-Maxwell} 仍然双曲。

在本构关系\eqref{align:ch03-本构eta} 下，\eqref{align:ch03-Maxwell} 与\eqref{align:ch03-高斯} 的任一光滑解还满足附加平衡律
\begin{align}
\partial_{t}P(B,D)+\operatorname{div}L(B,D)=-\rho E+B\wedge J,
\label{align:ch03-动量电磁}
\end{align}
其中
\begin{align}
P=B\wedge D,\qquad L=ED^{\top}+HB^{\top}+(\eta-E\cdot D-H\cdot B)I.
\label{align:ch03-PL}
\end{align}

特别地，在 Born--Infeld 情形下，
\begin{align}
P=Q,\qquad L=\eta^{-1}\left(I+BB^{\top}+DD^{\top}-QQ^{\top}\right).
\label{align:ch03-BI-PL}
\end{align}

读者应注意\eqref{align:ch03-电磁耗散} 与\eqref{align:ch03-动量电磁} 之间的差别：前者只依赖于\eqref{align:ch03-Maxwell}，而后者同时要求\eqref{align:ch03-Maxwell} 与\eqref{align:ch03-高斯} 成立。

沿方向 $\nu\in\mathbb{S}^{2}$ 以速度 $\sigma$ 传播的激波阵面，其 Rankine--Hugoniot 跳跃条件取为
\begin{align}
\begin{cases}
-\sigma[\![B]\!]=[\![E]\!]\wedge\nu\\
-\sigma[\![D]\!]=-[\![H]\!]\wedge\nu.
\end{cases}
\label{align:ch03-电磁RH}
\end{align}

为与\eqref{align:ch03-高斯} 相容，激波还应满足跳跃条件
\begin{align}
[\![B]\!]\cdot\nu=0,\qquad [\![D]\!]\cdot\nu=0.
\label{align:ch03-法向跳跃}
\end{align}

注意当 $\sigma\neq0$ 时\eqref{align:ch03-法向跳跃} 由\eqref{align:ch03-电磁RH} 推出。另一方面，\eqref{align:ch03-法向跳跃} 与双曲性一起排除了 $[\![E]\!]$、$[\![H]\!]$ 与 $\nu$ 共线、从而在 $\sigma=0$ 时满足\eqref{align:ch03-电磁RH} 的激波的可能性。

由\eqref{align:ch03-电磁耗散} 诱导出的\eqref{align:ch01-BV-激波不等式}%

不等式为
\begin{align}
-\sigma[\![\eta(B,D)]\!]+[\![E\wedge H]\!]\cdot\nu\leqslant0.
\label{align:ch03-电磁熵不等式}
\end{align}

还值得注意的是，在 Born--Infeld 本构关系\eqref{align:ch03-BornInfeld} 下，跳跃条件\eqref{align:ch03-电磁RH}、\eqref{align:ch03-法向跳跃} 推出
\begin{align}
-\sigma[\![P]\!]+[\![L]\!]\nu=0,
\label{align:ch03-电磁动量跳跃}
\end{align}
其中 $P$ 与 $L$ 由\eqref{align:ch03-BI-PL} 给出。这特别意味着：附加平衡律\eqref{align:ch03-动量电磁} 不仅对光滑解、甚至对\eqref{align:ch03-Maxwell} 与\eqref{align:ch03-高斯} 的 $BV$ 弱解都自动成立。

\subsection{磁流体力学}\label{subsection:ch03-MHD}

在电--机械现象的背景下会出现有趣的双曲平衡律方程组：连续介质热机械学的质量、动量、能量平衡律与 Maxwell 方程相耦合。作为说明性例子，我们在此考虑磁流体力学理论，它描述磁场与导电热弹性流体之间的相互作用。方程基于一系列简化假设，下面予以说明。

从 Maxwell 方程出发，电位移 $D$ 视为可忽略，于是由\eqref{align:ch03-Maxwell} 得 $J=\operatorname{curl}H$。磁感应 $B$ 与磁场 $H$ 通过经典关系 $B=\mu H$ 相联系。电场由流体在磁场中的运动产生，于是 $E=B\wedge v=\mu H\wedge v$。

流体是无热导热弹性介质，其热机械性质仍由本构关系\eqref{align:ch03-Euler本构} 描述。质量平衡\eqref{align:ch03-Euler}$_{1}$ 不受电磁场影响。另一方面，电磁场对流体施加力，应作为体积力计入动量平衡\eqref{align:ch03-Euler}$_{2}$。电场 $E$ 对该力的贡献视为可忽略，而磁场的贡献为 $J\wedge B=-\mu H\wedge\operatorname{curl}H$。利用恒等式
\begin{align}
-H\wedge\operatorname{curl}H=\operatorname{div}\left[HH^{\top}-\frac{1}{2}(H\cdot H)I\right],
\label{align:ch03-Maxwell应力}
\end{align}
该体积力可实现为 Maxwell 应力张量的散度。设无外体积力。为把电磁效应用入能量方程\eqref{align:ch03-Euler}$_{3}$，内能应加上电磁场能量 $\frac{1}{2}\mu|H|^{2}$，而 $\mu(H\wedge v)\wedge H=\mu|H|^{2}v-\mu(H\cdot v)H$（即 Poynting 向量）应加入通量。电磁能量产生 $-J\cdot E=-\mu(H\wedge v)\cdot\operatorname{curl}H$ 与电磁体积力的功率 $(J\wedge B)\cdot v=-\mu(H\wedge\operatorname{curl}H)\cdot v$ 相互抵消。

于是我们得到
\begin{align}
\begin{cases}
\partial_{t}\rho+\operatorname{div}(\rho v^{\top})=0\\
\partial_{t}(\rho v)+\operatorname{div}\left[\rho vv^{\top}-\mu HH^{\top}\right]+\operatorname{grad}\left[p(\rho,s)+\frac{1}{2}\mu|H|^{2}\right]=0\\
\partial_{t}\left[\rho\varepsilon(\rho,s)+\frac{1}{2}\rho|v|^{2}+\frac{1}{2}\mu|H|^{2}\right]\\
\quad+\operatorname{div}\left[\left(\rho\varepsilon(\rho,s)+\frac{1}{2}\rho|v|^{2}+p(\rho,s)+\mu|H|^{2}\right)v^{\top}-\mu(H\cdot v)H^{\top}\right]=\rho r\\
\partial_{t}H-\operatorname{curl}(v\wedge H)=0.
\end{cases}
\label{align:ch03-MHD系统}
\end{align}

上述平衡律方程组以 $U=(\rho,v,s,H)$ 为状态向量，当\eqref{align:ch03-Euler双曲条件} 成立时双曲。\eqref{align:ch03-MHD系统} 的热力学可容许解还应满足 $r=0$ 时的 Clausius--Duhem 不等式\eqref{align:ch03-Euler熵不等式}。由\eqref{align:ch03-Euler本构} 容易看出，\eqref{align:ch03-MHD系统} 的任一经典解都恒等地把\eqref{align:ch03-Euler熵不等式} 作为等式满足。于是 $-\rho s$ 是\eqref{align:ch03-MHD系统} 的熵，相应熵流为 $-\rho s v$。

方程组\eqref{align:ch03-MHD系统} 还须补充\eqref{align:ch03-高斯} 中相应的方程，即
\begin{align}
\operatorname{div}H=0.
\label{align:ch03-无散H}
\end{align}

磁流体力学比普通流体力学支持更丰富的激波族。与方程组\eqref{align:ch03-MHD系统} 相关联、沿方向 $\nu\in\mathbb{S}^{2}$ 以速度 $\sigma$ 传播的激波，其 Rankine--Hugoniot 跳跃条件为
\begin{align}
\begin{cases}
[\![\rho(v\cdot\nu-\sigma)]\!]=0\\
[\![\rho(v\cdot\nu-\sigma)v+\left(p+\frac{1}{2}\mu|H|^{2}\right)\nu-\mu(H\cdot\nu)H]\!]=0\\
[\![(v\cdot\nu-\sigma)\left(\rho\varepsilon+\frac{1}{2}\rho|v|^{2}+\frac{1}{2}\mu|H|^{2}\right)+(v\cdot\nu)\left(p+\frac{1}{2}\mu|H|^{2}\right)-\mu(H\cdot\nu)(H\cdot v)]\!]=0\\
[\![(v\cdot\nu-\sigma)H-(H\cdot\nu)v]\!]=0.
\end{cases}
\label{align:ch03-MHD-RH}
\end{align}

它们还须补充由\eqref{align:ch03-无散H} 诱导出的跳跃条件
\begin{align}
[\![H\cdot\nu]\!]=0,
\label{align:ch03-MHD法向跳跃}
\end{align}
它断言磁场的跃度必与激波面相切。

联立\eqref{align:ch03-MHD-RH} 与\eqref{align:ch03-MHD法向跳跃}，经过冗长计算可得\eqref{align:ch03-Euler-RH消速度} 的类比：
\begin{align}
\rho_{\pm}(v_{\pm}\cdot\nu-\sigma)\left\{[\![\varepsilon]\!]+\frac{1}{2}(p_{+}+p_{-})[\![\rho^{-1}]\!]-\frac{1}{4}[\![\rho^{-1}]\!][\![H]\!]^{2}\right\}=0.
\label{align:ch03-MHD-RH消速度}
\end{align}

跳跃条件\eqref{align:ch03-MHD-RH}、\eqref{align:ch03-MHD法向跳跃} 支持三类激波：
\begin{enumerate}
\item\label{item:ch03-MHD接触} \emph{接触间断}：类似于流体力学中的涡层，速度 $\sigma=v_{\pm}\cdot\nu$，且 $[\![p]\!]=0$、$[\![v]\!]=0$、$[\![H]\!]=0$，但 $[\![\rho]\!]\neq0$。这种跳跃不产生熵。
\item\label{item:ch03-MHD横波} \emph{横波激波}：$[\![\rho]\!]=0$、$[\![p]\!]=0$、$[\![s]\!]=0$。此时 $[\![v]\!]$ 与 $[\![H]\!]$ 必共线，
\begin{align}
[\![v]\!]=\pm\left(\frac{\mu}{\rho}\right)^{\frac{1}{2}}[\![H]\!],
\label{align:ch03-MHD横波关系}
\end{align}
特别地速度跳跃与激波面相切，即 $[\![v\cdot\nu]\!]=0$。此外，跨过这种激波磁场强度连续，$[\![|H|]\!]=0$。
\item\label{item:ch03-MHD压缩} \emph{压缩激波}：跨过它所有场 $\rho,p,\eta,v,H$ 都有非零跃度。当 Clausius--Duhem 不等式\eqref{align:ch03-Euler熵不等式} 满足时，这种激波称为压缩激波，因为密度、压力与熵跨过它都增加。压缩激波形成两个子类，按速度区分，分别称为快激波与慢激波。跨过快压缩激波磁场强度增强，跨过慢压缩激波磁场强度减弱。
\end{enumerate}

相应地存在各型弱磁流体波，形成复杂的几何图案。

\end{document}
